\begin{abstract}
Emerging studies, exemplified by BitNet~\cite{bitnet}, are ushering in an era of 1-bit Large Language Models (LLMs). In this paper, we propose a 1-bit LLM version, termed \textbf{\text{BitNet b1.58}}, characterized by representing every individual parameter (or weight) as a value in the ternary set \{-1, 0, 1\}. This model attains parity with full-precision (such as FP16 or BF16) Transformer LLMs of equivalent size and trained on identical token counts, excelling in both perplexity and downstream task benchmarks. Notably, it achieves substantial gains in latency, memory footprint, throughput, and energy efficiency. More fundamentally, the 1.58-bit LLM introduces a revised scaling law and a systematic approach for developing next-generation LLMs that combine state-of-the-art performance with reduced cost. It further paves the way for a novel computational paradigm and fosters the development of dedicated hardware tailored to the demands of 1-bit LLM architectures.
\end{abstract}

\section{The Era of 1-bit LLMs}

Artificial intelligence has witnessed a dramatic expansion in both the scale and capabilities of Large Language Models (LLMs) in recent years. These advances have driven significant improvements across various natural language processing tasks. However, as these models continue to grow, substantial obstacles have emerged related to practical deployment and environmental as well as financial concerns stemming from escalated energy usage.

One solution to these issues involves applying post-training quantization methods to compress model precision for inference~\cite{smoothquant, gptq, quip, quip_sharp}. By reducing the precision of model weights and activations, such approaches achieve notable reductions in memory and computational overhead for LLMs. The industry has been moving steadily from 16-bit precision toward more compact formats, including 4-bit representations~\cite{gptq, awq}. Nonetheless, despite its prevalence in real-world applications, post-training quantization generally fails to achieve optimal results.

Recent advancements in 1-bit neural architectures, exemplified by BitNet~\cite{bitnet}, offer a promising strategy for decreasing the operational costs of LLMs without sacrificing their effectiveness. Standard LLMs typically utilize 16-bit floating-point representations (FP16 or BF16), and most computational resources are allocated to matrix multiplication operations. As such, a significant portion of computation is devoted to floating-point addition and multiplication. Conversely, in BitNet, matrix multiplications rely solely on integer addition, which considerably reduces energy requirements in comparison to floating-point operations. Since power consumption is often the bottleneck for computational throughput in modern hardware, these reductions can also enable more rapid processing.

Beyond the computational savings, transferring model weights from off-chip DRAM to on-chip accelerator memory (such as SRAM) during inference incurs considerable overhead. While increasing SRAM capacity has been explored to alleviate this bottleneck, such solutions are typically much more costly than DRAM. One clear advantage of 1-bit LLMs is their dramatically reduced memory usage in both storage requirements and memory bandwidth, compared to models using full-precision weights. This substantial reduction facilitates faster and more cost-effective weight transfers, thus accelerating inference.

Here, we present a notable extension of 1-bit LLMs, denoted as \textbf{\text{BitNet b1.58}}, where all parameters are chosen from a ternary set \{-1, 0, 1\}. By incorporating a zero value to the binary weight set used in original BitNet, the effective bit representation becomes 1.58 bits. \text{BitNet b1.58} preserves the foundational advantages of the initial 1-bit BitNet, including its novel computation model that largely eliminates multiplications during matrix operations and allows for highly efficient implementation. It maintains the same power efficiency as its predecessor, while significantly outperforming FP16 LLMs in memory usage, throughput, and latency. Additionally, \text{BitNet b1.58} brings two notable enhancements. First, the inclusion of zero in the weights enables explicit feature selection, thereby strengthening the model’s representational power and yielding substantial gains for 1-bit LLM performance. Second, our empirical results indicate that \text{BitNet b1.58}, when trained under identical configurations (such as model size and training tokens) and for models at least 3B in size, achieves parity with FP16 models with respect to perplexity and end-task results.

\section{BitNet b1.58}

\text{BitNet b1.58} extends the BitNet framework, a Transformer variant that substitutes \emph{nn.Linear} layers with \emph{BitLinear} layers. This model is initialized and trained from scratch using weights quantized to 1.58 bits and activations quantized to 8 bits. Distinct from the original BitNet, several updates have been incorporated, which are outlined below.

\paragraph{Quantization Function.} We employ an \emph{absmean} quantization procedure to restrict weights to the set \{-1, 0, +1\}. Specifically, the weight matrix is normalized by its mean absolute value, after which each element is discretized to the nearest of -1, 0, or +1:

\begin{equation}
    \widetilde{W} = \mathrm{RoundClip}(\frac{W}{\gamma + \epsilon}, -1, 1),
\end{equation}

\begin{equation}
    \mathrm{RoundClip}(x, a, b) = \max(a, \min(b, \mathrm{round}(x))),
\end{equation}

\begin{equation}
    \gamma = \frac{1}{nm}\sum_{ij} |W_{ij}|.
\end{equation}

For activation quantization, we retain the procedure introduced in BitNet, with the exception that pre-nonlinearity activations are not rescaled to the interval $[0, Q_b]$. Rather, activations are linearly scaled for each token to $[-Q_b, Q_b]$, thereby omitting the use of zero-point quantization. This alteration not only simplifies implementation and system-level optimization but, based on our experiments, leads to an insubstantial effect on overall model performance.

\paragraph{LLaMA-alike Components.}

The LLaMA~\cite{llama, llama2} architecture has become the prevailing foundation for open-source LLM development. Aligning with this ecosystem, our implementation of \text{BitNet b1.58} adopts various LLaMA-inspired components, namely RMSNorm~\cite{rmsnorm}, SwiGLU~\cite{swiglu}, rotary embedding~\cite{rope}, and omits all biases from the model. Such choices facilitate effortless integration of \text{BitNet b1.58} with prominent open-source libraries and toolkits (for example, Huggingface, vLLM~\cite{vllm}, and llama.cpp\footnote{\url{https://github.com/ggerganov/llama.cpp}}).

\section{Results}

We conducted a comparative analysis of \text{BitNet b1.58} and a faithfully reproduced FP16 \text{LLaMA LLM} across multiple model scales. For parity in evaluation, both models underwent pre-training on the RedPajama dataset~\cite{redpajama} over 100 billion tokens. Zero-shot performance assessments were carried out on diverse language understanding benchmarks, namely ARC-Easy~\cite{arc}, ARC-Challenge~\cite{arc}, Hellaswag~\cite{hellaswag}, Winogrande~\cite{winoGrande}, PIQA~\cite{piqa}, OpenbookQA~\cite{openbookqa}, and BoolQ~\cite{boolq}. Additionally, validation perplexity scores were computed using WikiText2~\cite{wikitext2} and C4~\cite{c4} datasets.

We evaluated GPU memory usage and inference latency for both \text{LLaMA LLM} and \text{BitNet b1.58}, utilizing the FasterTransformer\footnote{\url{https://github.com/NVIDIA/FasterTransformer}} codebase, which provides highly efficient GPU-based inference for LLMs. For \text{BitNet b1.58}, the implementation incorporated the 2-bit kernel available through Ladder~\cite{ladder}. We focused on reporting the inference cost per output token, as this metric is most indicative of practical deployment efficiency.

Table~\ref{tab:ppl} presents both perplexity measurements and resource consumption statistics for \text{BitNet b1.58} compared to \text{LLaMA LLM}. Notably, starting at the 3B parameter model size, \text{BitNet b1.58} achieves perplexity on par with full precision \text{LLaMA LLM}, while attaining a 2.71$\times$ reduction in inference time and a 3.55$\times$ decrease in GPU memory usage. Of particular interest, \text{BitNet b1.58} at 3.9B parameters demonstrates 2.4$\times$ faster throughput and uses 3.32$\times$ less memory, while clearly outperforming the \text{LLaMA LLM} 3B model on accuracy.

Table~\ref{tab:zero-shot} details the results for zero-shot accuracy across all evaluated tasks. We adhered to the established \emph{lm-evaluation-harness}\footnote{\url{https://github.com/EleutherAI/lm-evaluation-harness}} methodology for these evaluations. The findings indicate that the performance discrepancy between \text{BitNet b1.58} and \text{LLaMA LLM} consistently diminishes as model size increases, with parity in performance reached at the 3B size. Echoing the perplexity observations, end-task evaluation demonstrates that \text{BitNet b1.58} at 3.9B not only consumes less memory and exhibits lower latency, but also surpasses \text{LLaMA LLM} 3B in accuracy. Collectively, these outcomes establish \text{BitNet b1.58} as a Pareto improvement relative to prevailing large language models.

\paragraph{Memory and Latency}

We expanded our experiments to include larger models at the 7B, 13B, and 70B scales and analyzed the associated computational costs. Figure~\ref{fig:latency-memory-energy} presents the relationship between model size and both latency and memory usage, indicating that efficiency gains become more pronounced as the models grow larger. Specifically, the \text{BitNet b1.58} 70B model achieves a 4.1-fold speed increase over the \text{LLaMA LLM} baseline. This acceleration is primarily attributed to the increasing time requirements of the \emph{nn.Linear} layer as model dimensions expand. A comparable trend emerges with respect to memory utilization, where the full-precision embedding accounts for a progressively smaller share in larger models. The experiments utilized a 2-bit kernel for both latency and memory assessment, suggesting that further optimization could still bring additional reductions in cost.

\paragraph{Energy}

Additionally, we quantified the energy consumption associated with arithmetic operations for both \text{BitNet b1.58} and \text{LLaMA LLM}. The primary focus was placed on matrix multiplication, as it represents the largest component of computational expense in LLMs. Figure~\ref{fig:energy} visualizes how energy costs are distributed. In \text{BitNet b1.58}, the majority of computation involves INT8 additions, whereas \text{LLaMA LLM} relies heavily on both FP16 additions and multiplications. Drawing on the energy model proposed in~\cite{energycost, pokebnn}, the matrix multiplication step in \text{BitNet b1.58} demonstrates a 71.4-fold reduction in arithmetic energy consumption compared to its FP16 counterpart on 7nm hardware. We also provide comprehensive measurements of end-to-end energy usage for inference involving 512 tokens. These findings reveal that with increasing model size, \text{BitNet b1.58} delivers ever greater energy efficiency relative to the FP16 \text{LLaMA LLM} baseline. This advantage arises as the proportion of computation devoted to \emph{nn.Linear} scales with model size, while other components’ relative cost diminishes for larger models.

\paragraph{Throughput}

We evaluate the throughput of \text{BitNet b1.58} and the 70B-parameter \text{LLaMA LLM} deployed on two A100 GPUs with 80GB memory each, employing pipeline parallelism~\cite{gpipe} to allow the execution of \text{LLaMA LLM} 70B within the given hardware constraints. The batch size was progressively increased until GPU memory was saturated, with each input sequence having a length of 512 tokens. According to Table~\ref{tab:throughput}, \text{BitNet b1.58} 70B accommodates up to 11-fold greater batch sizes compared to \text{LLaMA LLM}, culminating in an 8.9-fold increase in throughput.

\textbf{\text{BitNet b1.58} establishes a novel scaling law linking model quality and inference resource cost}. To provide context, the findings presented in Figures~\ref{fig:latency-memory-energy} and \ref{fig:energy} enable the following correspondences between 1.58-bit and 16-bit models of varying sizes.

\begin{itemize}
    \item The 13B BitNet b1.58 model demonstrates lower latency, memory footprint, and energy demand than the 3B FP16 LLM.
    \item The 30B BitNet b1.58 model surpasses the 7B FP16 LLM with respect to latency, memory usage, and energy efficiency.
    \item The 70B BitNet b1.58 model outperforms the 13B FP16 LLM in latency, memory consumption, and energy requirements.
\end{itemize}

\paragraph{Training with 2T Tokens}

The quantity of training tokens significantly impacts the effectiveness of LLMs. To assess the token-scaling characteristics of \text{BitNet b1.58}, we conducted training on a dataset of 2T tokens, adhering to the StableLM-3B~\cite{StableLM-3B-4E1T} data preparation pipeline, representative of leading open-source 3B models. Evaluation of both models was performed on a suite of benchmarks: Winogrande~\cite{winoGrande}, PIQA~\cite{piqa}, SciQ~\cite{sciq}, LAMBADA~\cite{lambada}, and ARC-easy~\cite{arc}. Zero-shot accuracy metrics are summarized in Table~\ref{tab:2t}. For tasks involving both accuracy and normalized accuracy, their mean is reported. Results for StableLM 3b at 2T are reproduced from its technical documentation. The experimental evidence demonstrates that \text{BitNet b1.58} achieves higher accuracy across all evaluated tasks, highlighting the robust generalization capacity of 1.58-bit LLMs.

\section{Discussion and Future Work}

\textbf{1-bit Mixture-of-Experts (MoE) LLMs}

Mixture-of-Experts (MoE) architectures have demonstrated effectiveness as a resource-efficient solution for large language models. Although they notably reduce computational FLOPs, memory usage and inter-chip communication still pose significant obstacles to their practical deployment. Utilizing 1.58-bit LLMs offers promising avenues to mitigate these issues. With a lower memory footprint, the demand for hardware resources to deploy MoE models diminishes, and the latency associated with transmitting activations across different devices is markedly lessened. In an optimal scenario where the full model resides on a single chip, communication overheads can be virtually eliminated.

\textbf{Native Support of Long Sequence in LLMs}

The requirement to process extended sequences has become increasingly important in contemporary LLM applications. However, memory constraints arising from key-value (KV) caches present a considerable barrier to efficient long sequence inference. The development of \text{BitNet b1.58} marks progress in this area by halving activation storage requirements from 16 bits to 8 bits, thereby permitting twice the context length with existing resources. There is potential for additional compression—potentially reducing activations to 4 bits or even lower in 1.58-bit LLMs—which we propose to explore in subsequent research.

\textbf{LLMs on Edge and Mobile}

Deploying 1.58-bit LLMs could transform the landscape of language model capabilities on edge and mobile platforms. Memory and computational constraints often limit the feasible scale and performance of LLMs in such environments. The memory efficiency and reduced energy requirements of 1.58-bit LLMs make it possible to bring sophisticated LLM applications to these resource-constrained devices. This advancement expands the potential application domains and enhances device utility. Additionally, 1.58-bit LLMs offer improved compatibility with CPUs—the primary processors in mobile and edge hardware—which means that models like \text{BitNet b1.58} can execute more efficiently and effectively on these platforms, boosting both performance and utility.

\textbf{New Hardware for 1-bit LLMs}

Recent innovations such as Groq\footnote{\url{https://groq.com/}} have highlighted the possibilities inherent in dedicated hardware (e.g., LPUs) tailored for LLM computation. To capitalize on the computational patterns introduced by BitNet~\cite{bitnet}, we propose further exploration and dedicated design of hardware and systems customized specifically for 1-bit LLMs, opening avenues for new computational frameworks and optimized deployment.

\end{document}